\SourceLeaf{273}%
\noindent Wählte man hingegen nach den obigen Dimensionen den zehnmillionsten Teil des Erdquadranten selbst als Einheit, so ergäbe sich
\[
\log A=6{,}8049902365
\]
\par\medskip\Needspace{4\baselineskip}\pdfbookmark[1]{Art. 6}{geodaesie1-art-6}\centerline{\textbf{6}}\nobreak\smallskip
\indent Die Breite auf der Kugel lässt sich ohne Weiteres nach Formel 3 aus der Breite auf dem Ellipsoid berechnen, wenn die Rechnung nur für wenige Fälle erforderlich ist. Für umfangreichere Anwendungen ist es dagegen vorteilhaft, eine Reihe zu verwenden. Die für ihre Entwicklung nötigen Formeln sollen hier angegeben werden.

Ich bezeichne eine unbestimmte Breite auf dem Ellipsoid, also einen unbestimmten Wert von \(90^0-w\), mit \(P+p\) und die entsprechende Breite auf der Kugel, also den Wert von \(90^0-U\), mit \(Q+q\). Nach dem Satz von \textsc{Taylor} gilt
\[
q=\frac{\mathrm{d}U}{\mathrm{d}w}\mathord{.}p-\frac{1}{2}\mathord{.}\frac{\mathrm{dd}U}{\mathrm{d}w^2}\mathord{.}pp+\frac{1}{6}\mathord{.}\frac{\mathrm{d}^3U}{\mathrm{d}w^3}\mathord{.}p^3-\frac{1}{24}\mathord{.}\frac{\mathrm{d}^4U}{\mathrm{d}w^4}\mathord{.}p^4+\text{und so weiter}
\]
wobei für die Ableitungen diejenigen bestimmten Werte einzusetzen sind, die zu \(p=0\) oder zu \(w=90^0-P\), \(U=90^0-Q\) gehören. Entwickelt man die Ableitungen der Reihe nach in ihrer allgemeinen Form, so erhält man
\[
\frac{\mathrm{d}U}{\mathrm{d}w}=\frac{\alpha(1-ee)\sin U}{(1-ee\cos w^2)\sin w}
\]
\[
\begin{aligned}
\frac{\mathrm{dd}U}{\mathrm{d}w^2}={}&\frac{\alpha(1-ee)\sin U}{(1-ee\cos w^2)^2\sin w^2}\\
&\quad\{\alpha(1-ee)\cos U-\cos w+ee(\cos w^3-2\cos w\sin w^2)\}
\end{aligned}
\]
\[
\begin{aligned}
\frac{\mathrm{d}^3U}{\mathrm{d}w^3}={}&\frac{\alpha(1-ee)\sin U}{(1-ee\cos w^2)^3\sin w^3}\{\alpha\alpha(1-ee)^2(\cos U^2-\sin U^2)\\
&-3\alpha(1-ee)\cos U(\cos w-ee(\cos w^3-2\cos w\sin w^2))\\
&+2\cos w^2+\sin w^2-ee(4\cos w^4-2\sin w^4)\\
&+e^4(2\cos w^6-\cos w^4\sin w^2+6\cos w^2\sin w^4)\}
\end{aligned}
\]
Die beiden folgenden Ableitungen habe ich ebenfalls entwickelt. Um Platz zu sparen, gebe ich ihre allgemeinen Formen hier nicht an.

\Needspace{7\baselineskip}
Setzt man \(w=90^0-P\), \(U=90^0-Q\) ein und ersetzt zugleich
\par\noindent \(\alpha\sin Q\) durch den Wert \(\sin P\) (nach Gleichung 10) und
\par\noindent \(\alpha\cos Q\) durch den Wert \(\frac{\cos P}{\cos\zeta\cos\eta}=\frac{\cos\theta\cos P}{\cos\varphi}\) (nach den Gleichungen 18, 16, 23)
\par\noindent und schreibt man abkürzend \(\cos P=c\), \(\sin P=s\), so erhält man
